% Bibliografía IEEE en orden de primera cita (orden de primera cita (IEEE)).
\begin{thebibliography}{34}
\bibitem{r1}
World Health Organization, \emph{WHO Recommendations for Care of the Preterm or Low-Birth-Weight Infant}. Geneva: WHO, 2022.

\bibitem{r2}
N. Charpak, R. Tessier, J. G. Ruiz et al., ``Twenty-year follow-up of kangaroo mother care versus traditional care,'' \emph{Pediatrics}, vol. 139, no. 1, e20162063, 2017.

\bibitem{r3}
N. Charpak and A. Montealegre-Pomar, ``Follow-up of Kangaroo Mother Care programmes in the last 28 years: results from a cohort of 57\,154 low-birth-weight infants in Colombia,'' \emph{BMJ Global Health}, vol. 8, no. 5, 2023.

\bibitem{r4}
Fundación Canguro, ``Modelo conceptual SEVARTI y diccionario de variables del Observatorio Canguro,'' especificación y material gráfico institucional proporcionados para este estudio, 2026 (material técnico no publicado).

\bibitem{r5}
S. E. Melo, O. O. Melo, and N. Suarez, ``Modelo multinivel de tasa global de ganancia de peso en el programa madre canguro en Bogotá,'' \emph{Revista de Salud Pública}, vol. 18, no. 2, pp. 238--250, 2016, doi: 10.15446/rsap.v18n2.44661.

\bibitem{r6}
Z. Ren, W. Gao, Q. Wang et al., ``Predictive role of NICU-related stress, postpartum depression trajectory and family coping on growth trajectory of moderate-to-late preterm infants: A longitudinal study,'' \emph{Journal of Advanced Nursing}, vol. 80, no. 8, pp. 3167--3178, 2024, doi: 10.1111/jan.16068.

\bibitem{r7}
S. Nel, U. D. Feucht, T. Botha, M. Arashi, and F. A. M. Wenhold, ``First-year growth patterns of preterm infants receiving kangaroo mother care: associations with early life factors and 1-year anthropometry,'' \emph{European Journal of Clinical Nutrition}, vol. 80, no. 1, pp. 62--72, 2026, doi: 10.1038/s41430-025-01662-6.

\bibitem{r8}
S. Baker and Y. Kandasamy, ``Machine learning for understanding and predicting neurodevelopmental outcomes in premature infants: a systematic review,'' \emph{Pediatric Research}, vol. 93, pp. 293--299, 2023.

\bibitem{r9}
G. M. Demirci et al., ``Predicting mental and psychomotor delay in very pre-term infants using machine learning,'' \emph{Pediatric Research}, vol. 95, pp. 668--678, 2024.

\bibitem{r10}
K. Carlton, J. Zhang, E. Cabacungan et al., ``Machine learning risk stratification for high-risk infant follow-up of term and late preterm infants,'' \emph{Pediatric Research}, vol. 99, pp. 1827--1835, 2026, doi: 10.1038/s41390-024-03338-6.

\bibitem{r11}
L.-Y. Wang, L.-Y. Wang, M.-I. Sung et al., ``Implementing explainable machine learning models for practical prediction of early neonatal hypoglycemia,'' \emph{Diagnostics}, vol. 14, no. 14, art. 1571, 2024, doi: 10.3390/diagnostics14141571.

\bibitem{r12}
M. G. Kahn, T. J. Callahan, J. Barnard et al., ``A harmonized data quality assessment terminology and framework for the secondary use of electronic health record data,'' eGEMs, vol. 4, no. 1, art. 1244, 2016, doi: 10.13063/2327-9214.1244.

\bibitem{r13}
N. G. Weiskopf and C. Weng, ``Methods and dimensions of electronic health record data quality assessment: enabling reuse for clinical research,'' \emph{Journal of the American Medical Informatics Association}, vol. 20, no. 1, pp. 144--151, 2013.

\bibitem{r14}
V. Chandola, A. Banerjee, and V. Kumar, ``Anomaly detection: A survey,'' \emph{ACM Computing Surveys}, vol. 41, no. 3, art. 15, 2009.

\bibitem{r15}
R. Koenker and G. Bassett, ``Regression quantiles,'' \emph{Econometrica}, vol. 46, no. 1, pp. 33--50, 1978.

\bibitem{r16}
Y. Romano, E. Patterson, and E. Candès, ``Conformalized quantile regression,'' in \emph{Advances in Neural Information Processing Systems}, vol. 32, 2019.

\bibitem{r17}
A. N. Angelopoulos and S. Bates, ``Conformal prediction: A gentle introduction,'' \emph{Foundations and Trends in Machine Learning}, vol. 16, no. 4, pp. 494--591, 2023.

\bibitem{r18}
B. Van Calster, D. J. McLernon, M. van Smeden, L. Wynants, and E. W. Steyerberg, ``Calibration: the Achilles heel of predictive analytics,'' \emph{BMC Medicine}, vol. 17, art. 230, 2019.

\bibitem{r19}
S. E. Davis, T. A. Lasko, G. Chen, and M. E. Matheny, ``Calibration drift among regression and machine learning models for hospital mortality,'' in \emph{AMIA Annual Symposium Proceedings}, vol. 2017, pp. 625--634, 2017.

\bibitem{r20}
A. Wong, E. Otles, J. P. Donnelly et al., ``External validation of a widely implemented proprietary sepsis prediction model in hospitalized patients,'' \emph{JAMA Internal Medicine}, vol. 181, no. 8, pp. 1065--1070, 2021, doi: 10.1001/jamainternmed.2021.2626.

\bibitem{r21}
S. M. Lundberg, G. Erion, H. Chen et al., ``From local explanations to global understanding with explainable AI for trees,'' \emph{Nature Machine Intelligence}, vol. 2, no. 1, pp. 56--67, 2020.

\bibitem{r22}
A. M. Salih, G. Menegaz, T. Pillay, and E. M. Boyle, ``Explainable artificial intelligence in paediatric: challenges for the future,'' \emph{Health Science Reports}, vol. 7, no. 12, e70271, 2024, doi: 10.1002/hsr2.70271.

\bibitem{r23}
M. Ghassemi, L. Oakden-Rayner, and A. L. Beam, ``The false hope of current approaches to explainable artificial intelligence in health care,'' \emph{The Lancet Digital Health}, vol. 3, no. 11, pp. e745--e750, 2021.

\bibitem{r24}
G. S. Collins, K. G. M. Moons, P. Dhiman et al., ``TRIPOD+AI statement: updated guidance for reporting clinical prediction models that use regression or machine learning methods,'' \emph{BMJ}, vol. 385, e078378, 2024, doi: 10.1136/bmj-2023-078378.

\bibitem{r25}
E. S. Page, ``Continuous inspection schemes,'' \emph{Biometrika}, vol. 41, no. 1/2, pp. 100--115, 1954.

\bibitem{r26}
G. Ke, Q. Meng, T. Finley et al., ``LightGBM: A highly efficient gradient boosting decision tree,'' in \emph{Advances in Neural Information Processing Systems}, vol. 30, pp. 3146--3154, 2017.

\bibitem{r27}
C. Leys, C. Ley, O. Klein, P. Bernard, and L. Licata, ``Detecting outliers: Do not use standard deviation around the mean, use absolute deviation around the median,'' \emph{Journal of Experimental Social Psychology}, vol. 49, no. 4, pp. 764--766, 2013, doi: 10.1016/j.jesp.2013.03.01

\bibitem{r28}
B. Iglewicz and D. C. Hoaglin, \emph{How to Detect and Handle Outliers}. Milwaukee, WI, USA: ASQC Quality Press, 1993.

\bibitem{r29}
World Health Organization, \emph{WHO Child Growth Standards: Length/Height-for-Age, Weight-for-Age, Weight-for-Length, Weight-for-Height and Body Mass Index-for-Age: Methods and Development}. Geneva: WHO, 2006.

\bibitem{r30}
World Health Organization and UNICEF, \emph{Recommendations for Data Collection, Analysis and Reporting on Anthropometric Indicators in Children under 5 Years Old}. Geneva: WHO, 2019.

\bibitem{r31}
D. Kreuzberger, N. Kühl, and S. Hirschl, ``Machine learning operations (MLOps): Overview, definition, and architecture,'' \emph{IEEE Access}, vol. 11, pp. 31866--31879, 2023, doi: 10.1109/ACCESS.2023.3262138.

\bibitem{r32}
M. Zaharia, A. Chen, A. Davidson et al., ``Accelerating the machine learning lifecycle with MLflow,'' \emph{IEEE Data Engineering Bulletin}, vol. 41, no. 4, pp. 39--45, 2018.

\bibitem{r33}
T. R. Fenton and J. H. Kim, ``A systematic review and meta-analysis to revise the Fenton growth chart for preterm infants,'' \emph{BMC Pediatrics}, vol. 13, art. 59, 2013.

\bibitem{r34}
J. Villar, F. Giuliani, Z. A. Bhutta et al., ``Postnatal growth standards for preterm infants: the Preterm Postnatal Follow-up Study of the INTERGROWTH-21st Project,'' \emph{The Lancet Global Health}, vol. 3, no. 11, pp. e681--e691, 2015, doi: 10.1016/S2214-109X(15)00163-1.

\end{thebibliography}
